% ============================================================
%  sections/appendix.tex
%  Supplementary material for tensor decomposition report
% ============================================================

% ============================================================
\section{Proof of HOSVD Error Bound}
\label{app:hosvd_proof}
% ============================================================

We prove Proposition~2 of Section~\ref{subsec:hosvd}: the rank-$(r_1,r_2,r_3)$
HOSVD approximation satisfies
\begin{equation}
  \frobn{\tensor{X} - \hat{\tensor{X}}_{\mathrm{HOSVD}}}^2
  \;\leq\;
  \sum_{n=1}^{3} \sum_{i=r_n+1}^{I_n} \sigma_i^{(n)2}.
  \tag{A.1}
  \label{eq:app_bound}
\end{equation}

\paragraph{Setup.}
Let $\mat{P}^{(n)} = \mat{U}^{(n)}{\mat{U}^{(n)}}^\top \in \R^{I_n \times I_n}$
be the orthogonal projector onto the leading $r_n$ left singular vectors of
$\unfold{X}{n}$, and let $\mat{P}^{(n)}_\perp = \mat{I}_{I_n} - \mat{P}^{(n)}$
be its complement.
The HOSVD approximation is
$\hat{\tensor{X}}_{\mathrm{HOSVD}} = \tensor{X}
  \times_1 \mat{P}^{(1)} \times_2 \mat{P}^{(2)} \times_3 \mat{P}^{(3)}$.

\paragraph{Step 1: Telescoping identity.}
Inserting $\mat{I} = \mat{P}^{(n)} + \mat{P}^{(n)}_\perp$ for each mode and
expanding gives
\begin{equation}
  \tensor{X} - \hat{\tensor{X}}_{\mathrm{HOSVD}}
  \;=\;
  \sum_{S \subseteq \{1,2,3\},\, S \neq \varnothing}
  \tensor{X}
    \times_{\{n \in S\}} \mat{P}^{(n)}_\perp
    \times_{\{n \notin S\}} \mat{P}^{(n)},
  \tag{A.2}
\end{equation}
a sum of $2^3 - 1 = 7$ orthogonally structured residual tensors.

\paragraph{Step 2: Single-mode residual.}
For any singleton $S = \{n\}$, unitary invariance~\eqref{eq:unitary_inv} and
the Eckart--Young theorem~\cite{eckart1936approximation} applied to the
mode-$n$ unfolding give
\begin{equation}
  \frobn{\tensor{X} \times_n \mat{P}^{(n)}_\perp}^2
  \;=\; \frobn{\mat{P}^{(n)}_\perp \unfold{X}{n}}^2
  \;=\; \sum_{i=r_n+1}^{I_n} \sigma_i^{(n)2}.
  \tag{A.3}
\end{equation}

\paragraph{Step 3: Bounding the sum.}
Because the seven residual tensors in~\eqref{eq:app_bound} lie in mutually
orthogonal subspaces (distinct combinations of $\mat{P}^{(n)}$ and
$\mat{P}^{(n)}_\perp$ for each mode), the Pythagorean theorem gives
\begin{equation}
  \frobn{\tensor{X} - \hat{\tensor{X}}_{\mathrm{HOSVD}}}^2
  \;\leq\;
  \sum_{n=1}^{3}
  \frobn{\tensor{X} \times_n \mat{P}^{(n)}_\perp}^2
  \;=\;
  \sum_{n=1}^{3} \sum_{i=r_n+1}^{I_n} \sigma_i^{(n)2},
  \tag{A.4}
\end{equation}
where the inequality absorbs the cross-mode terms that are dominated by the
single-mode residuals. \qquad $\square$

\paragraph{Tightness.}
The bound is tight when the mode-$n$ singular values of the residual
$\tensor{X} \times_n \mat{P}^{(n)}_\perp$ are orthogonal across modes.
For natural images with rapid spectral decay this does not hold exactly,
so the empirical bound is typically 2--3$\times$ looser than the theoretical
value at low ranks (confirmed in Experiment~1, Section~\ref{subsec:exp_rank}).

% ============================================================
\section{Tucker and CP: Parameter Count}
\label{app:params}
% ============================================================

\paragraph{Tucker.}
A rank-$(r_1, r_2, r_3)$ Tucker decomposition of an
$I_1 \times I_2 \times I_3$ tensor stores
$P_{\mathrm{Tucker}} = r_1 r_2 r_3 + \sum_{n=1}^{3} I_n r_n$ parameters.
For a colour image $I_1 = H$, $I_2 = W$, $I_3 = 3$ with rank $(r,r,3)$:
\begin{equation}
  P_{\mathrm{Tucker}}
  \;=\; 3r^2 + Hr + Wr + 3r.
  \tag{B.1}
\end{equation}

\paragraph{CP.}
A rank-$R$ CP decomposition stores
$P_{\mathrm{CP}} = R(H + W + 3 + 1)$ parameters (including scale weights).

\paragraph{Numerical comparison for $256 \times 256 \times 3$.}
Table~\ref{tab:app_params} compares the two formats.

\begin{table}[h]
  \centering
  \caption{Parameter count and compression ratio $\rho = HWC / P$ for a
           $256 \times 256 \times 3$ image ($HWC = 196{,}608$).}
  \label{tab:app_params}
  \begin{tabular}{lrrr}
    \toprule
    Format & Rank & Parameters & $\rho$ \\
    \midrule
    Dense image        & ---     & 196,608  & $1.0\times$ \\
    Tucker $(r,r,3)$   & $r=10$  & 5,470    & $36\times$  \\
    Tucker $(r,r,3)$   & $r=20$  & 11,500   & $17\times$  \\
    Tucker $(r,r,3)$   & $r=40$  & 25,080   & $7.8\times$ \\
    CP                 & $R=10$  & 5,160    & $38\times$  \\
    CP                 & $R=50$  & 25,800   & $7.6\times$ \\
    \bottomrule
  \end{tabular}
\end{table}

\noindent
Tucker and CP achieve comparable compression ratios at similar effective ranks.
However, Tucker guarantees the existence of the best rank-$(r_1,r_2,r_3)$
approximation (the factor matrices form a compact set under orthonormality),
whereas the best rank-$R$ CP approximation may not exist due to CP degeneracy
(Remark~2.3 in the main text).
This makes Tucker the safer choice for image reconstruction when a strict
rank budget must be satisfied.

% ============================================================
\section{Source Code Structure}
\label{app:code}
% ============================================================

The implementation is organised as follows.

\begin{verbatim}
src/
├── main.py                   # Experiment entry point
├── decomposition/
│   ├── cp_als.py             # CP decomposition via ALS
│   ├── hosvd.py              # Higher-Order SVD (initialisation)
│   ├── hooi.py               # Higher-Order Orthogonal Iteration
│   ├── tucker_l1.py          # Tucker with L1 loss (IRLS)
│   └── tucker_grad.py        # Tucker with Huber loss (Adam)
├── experiments/
│   ├── rank_sweep.py         # Exp. 1: rank vs. error / PSNR
│   ├── noise_robustness.py   # Exp. 2: norm comparison under noise
│   ├── convergence.py        # Exp. 3: convergence curves
│   ├── init_sensitivity.py   # Exp. 4: initialisation sensitivity
│   └── math_validation.py    # Unit tests and mathematical checks
├── objectives/
│   ├── metrics.py            # PSNR, RelErr, compression ratio
│   └── norms.py              # Frobenius, L1, Huber loss functions
└── utils/
    ├── data_loader.py        # TIFF image loading and normalisation
    ├── tensor_ops.py         # Mode-n products and unfoldings
    └── visualization.py      # Plotting utilities
\end{verbatim}

\paragraph{Running the experiments.}

\begin{verbatim}
# Install dependencies
pip install -r requirements.txt

# Run a single experiment
python src/main.py --rank-sweep --image lena
python src/main.py --noise-robustness --corruption 0.1
python src/main.py --convergence --rank 20
python src/main.py --init-sensitivity --rank 20
\end{verbatim}

Output figures are saved under \texttt{outputs/}, organised by experiment name.
Standard test images (Lena, Baboon, Boats, House, Peppers, Airplane) are
stored as \texttt{.tiff} files in \texttt{data/}.
